\section{Purpose and Document Ecosystem}

This guide exists because effective project advising requires more than technical expertise---it requires understanding of the course structure, grading philosophy, common failure modes, and when to intervene versus when to let students struggle productively. Whether you are a first-time PA or a veteran, this document provides a shared baseline for how we run Capstone Design.

\subsection{The Four-Document Ecosystem}

The Capstone Design program is supported by four coordinated documents. Each serves a distinct audience and purpose:

\begin{enumerate}[nosep,leftmargin=1.5em]
\item \textbf{Course Text} (22 chapters, 5 parts) --- The comprehensive reference for both students and PAs. Covers engineering design methodology, project management, professional skills, safety, and compliance. When students ask ``why do we have to do this?'' the Course Text provides the pedagogical justification and worked examples.

\item \textbf{Student Guide} --- A sprint-by-sprint navigation aid for students. Tells them what to focus on each sprint, common pitfalls to avoid, and deliverable checklists. Written in a direct, student-facing tone.

\item \textbf{Canvas} --- The authoritative source for assignments, rubrics, deadlines, and the Course Syllabus. All grading policies, late-work rules, and academic integrity standards live here. The Course Syllabus is the contractual document.

\item \textbf{PA Guide} (this document) --- Instructor-facing reference covering facilitation, grading calibration, team management, and design review scripts. Not intended for student distribution.
\end{enumerate}

\begin{tipbox}
When a student asks a question, point them to the right document. Course mechanics and deadlines: Canvas. Design methodology and theory: Course Text. Sprint-level priorities: Student Guide. Resist the urge to re-explain what is already documented---teach students to use the resources.
\end{tipbox}

\subsection{Quick Start: Your First Week as a PA}

If you are new to the PA role, complete the following during your first week:

\begin{enumerate}[nosep,leftmargin=1.5em]
\item \textbf{Read the Course Syllabus on Canvas.} Understand grading weights, attendance policies, and academic integrity expectations. You will be asked about these.

\item \textbf{Skim the Student Guide.} Understand what students are being told. This prevents contradictory messaging.

\item \textbf{Read Sections~2--4 of this guide.} These cover your role, the sprint structure, and design review facilitation---the core of your weekly work.

\item \textbf{Review your team assignments.} Note the project type (Build, Paper, Competition, or Hybrid), the client organization, and any special constraints (e.g., ITAR, competition deadlines).

\item \textbf{Attend the PA orientation meeting.} The CF and CLT Director will brief you on this semester's schedule, any policy changes, and administrative procedures.

\item \textbf{Set up your meeting cadence.} Schedule a recurring weekly or biweekly meeting with each of your teams. Consistency matters more than frequency.

\item \textbf{Introduce yourself to your teams' clients.} A brief email establishing that you are the PA and providing your contact information sets the right tone.
\end{enumerate}

\begin{paaction}
Before the first sprint review, review the Course Text chapters covering the current sprint's topics. You do not need to read all 22 chapters---but you should always be at least one sprint ahead of your teams in terms of content familiarity.
\end{paaction}


% ═══════════════════════════════════════
